\documentclass[10pt,a4paper]{amsart}
\usepackage[margin=19mm]{geometry}
\usepackage{amsmath,amssymb,mathtools}
\usepackage[scr=rsfs,cal=euler]{mathalfa}
\usepackage{xfrac}
\usepackage[unicode,hidelinks,bookmarks=false]{hyperref}
\newcommand{\ie}{i.e.\ }
\newcommand{\cf}{cf.\ }
\newcommand{\resp}{respectively\ }
\newcommand{\Subsection}{\subsection}
\providecommand{\nobreakdash}{\nobreak\hbox{-}}
\setlength{\emergencystretch}{2em}
\multlinegap0pt
\usepackage{amssymb}
\usepackage{mathtools}
\usepackage[shortlabels]{enumitem}
\usepackage{xcolor}
\newtheorem{theorem}{Theorem}
\usepackage{cleveref}
\crefname{theorem}{Theorem}{Theorems}
\newcommand{\E}{E}
\newcommand{\Gs}{(G,\sigma)}
\newcommand{\V}{V}
\title{\bfseries Every signed planar graph is $5$-choosable:\\ A short proof and refinements}
\author{Pie Desire Ebode Atangana, Maxwell Ndognkon Manga}
\date{Source: arXiv:2605.22860v1 (CC BY 4.0)}
\begin{document}
\maketitle

\noindent\emph{Source and licence.} \bfseries Every signed planar graph is $5$-choosable:\\ A short proof and refinements. Version arXiv:2605.22860v1 (CC BY 4.0), distributed by arXiv under the Creative Commons Attribution 4.0 International licence, http://creativecommons.org/licenses/by/4.0/, as recorded in the arXiv licence metadata for this exact version and shown on the abstract page https://arxiv.org/abs/2605.22860v1. Full licence notice: \texttt{licenses/arxiv-2605.22860-CC-BY-4.0.txt}.\par\smallskip

\noindent\textit{Editorial scope.} Counted unit: the article's theorem thm:main (source0.tex lines 307--322). The article's complete original proof of it is delivered with it (source0.tex lines 663--790, one \texttt{proof} environment of 128 lines, which the article places away from the statement). No mathematics is written, repaired or re-derived: the statement and the proof are the article's own lines, in the article's own notation, and no retained line is substituted or re-flowed.  No further source-local context is needed: every cross-reference of every retained line resolves inside the two delivered spans.  Nothing was omitted from any retained span, so the package carries no omission record.  No retained line contains a citation, so the wrapper carries no bibliography and no bibliography entry.  No retained line contains a figure environment, an image-inclusion command or drawing code, so no image is delivered and the package declares no asset dependency.  Editorial formatting only, in addition: an AMS \texttt{amsart} preamble that restates the article's own declarations and macros used by the retained lines, namely the environments \texttt{theorem} on separate counters, together with the standard packages those lines need.  The retained environments print in this document as Theorem 1 in order of appearance, whereas the article prints the same items with the numbering of its own full document.  The complete native source of the article as distributed in the arXiv e-print for this exact version is bundled unchanged as \texttt{sources/201164/source0.tex}.
\begin{theorem}[Main Extension Theorem]\label{thm:main}
Let $\Gs$ be a signed near-triangulation with outer cycle
$C\colon v_1v_2\cdots v_pv_1$ ($p\ge 3$), and let $L$ be a list
assignment satisfying the following three hypotheses. First, the
two adjacent boundary vertices $v_1$ and $v_2$ are precolored with
values $c_1\in L(v_1)$ and $c_2\in L(v_2)$ in such a way that the
proper-coloring condition is already satisfied on the edge $v_1v_2$,
that is, $c_1\ne\sigma(v_1v_2)\,c_2$; we shall refer to this as
condition~(i). Second, the remaining boundary vertices each carry
a list of size at least three, so that $|L(v)|\ge 3$ for every
$v\in\V(C)\setminus\{v_1,v_2\}$, which we call condition~(ii).
Third, every interior vertex carries a list of size at least five,
so that $|L(v)|\ge 5$ for every $v\in\V(G)\setminus\V(C)$, which
we call condition~(iii). Then the precoloring of $\{v_1,v_2\}$
extends to a proper $L$-coloring of $\Gs$.
\end{theorem}

\begin{proof}[Proof of \cref{thm:main}]
We use induction on $n:=|\V(G)|$.

In the base case, $n=3$, so $G=C=v_1v_2v_3$. Vertices $v_1,v_2$ are
precolored consistently with~(i), and the constraints on $c(v_3)$
are
\[
c(v_3)\;\ne\;\sigma(v_1v_3)\,c_1
\quad\text{and}\quad
c(v_3)\;\ne\;\sigma(v_2v_3)\,c_2,
\]
that is, at most two forbidden values. Since $|L(v_3)|\ge 3$
by~(ii), a valid color is available.

For the inductive step, $n\ge 4$, we split on whether $C$ has a
chord. The two cases are quite different in flavor: a chord allows
us to recurse on two strictly smaller pieces, while a chordless
boundary forces us to remove a single boundary vertex.

Consider first the case in which $C$ has a chord. Let
$v_iv_j\in\E(G)$ be a chord with $1\le i<j\le p$, $j-i\ge 2$, and
$\{i,j\}\ne\{1,p\}$. The chord $v_iv_j$ together with $C$ splits
the disk bounded by $C$ into two regions; correspondingly $G$
decomposes into two near-triangulations $G_1$ and $G_2$, sharing
the edge $v_iv_j$. Relabeling if necessary, we may assume that
both $v_1$ and $v_2$ belong to $G_2$, with the arc
$v_1\cdots v_iv_jv_p\cdots v_2$ on its boundary.

We apply the induction hypothesis to $\Gs|_{G_2}$ with the
precoloring of $v_1,v_2$. The hypotheses are inherited, since
$L|_{V(G_2)}$ has size at least $3$ on
$\V(\partial G_2)\setminus\{v_1,v_2\}$ and size at least $5$ in the
interior. The induction yields a proper $L$-coloring $c_2$ of
$\Gs|_{G_2}$, which in particular determines values
$b_i:=c_2(v_i)\in L(v_i)$ and $b_j:=c_2(v_j)\in L(v_j)$ with
$b_i\ne\sigma(v_iv_j)\,b_j$, since the chord $v_iv_j\in\E(G_2)$ is
properly colored. We then turn to $\Gs|_{G_1}$ with the modified
list assignment
\[
   L'(v_i)=\{b_i\},\quad L'(v_j)=\{b_j\},\quad
   L'(v)=L(v)\text{ otherwise.}
\]
The outer cycle of $G_1$ is the cycle whose two precolored adjacent
vertices are $v_i$ and $v_j$, playing the roles of $v_1,v_2$, so
the inductive hypotheses are again met: the remaining boundary
lists have size at least $3$, and interior vertices of $G_1$ are
interior in $G$ with lists of size at least $5$. Induction yields a
proper $L'$-coloring $c_1$ of $\Gs|_{G_1}$, and since $c_1$ and
$c_2$ agree on the chord $v_iv_j$, their union is a proper
$L$-coloring of $\Gs$. This concludes the chord case.

Consider next the case in which $C$ has no chord. Here we
eliminate the boundary vertex $v_p$. Listing the neighbors of $v_p$
in $G$ in clockwise order around $v_p$ in the planar embedding
gives a sequence
\[
v_1,\,u_1,\,u_2,\,\dots,\,u_m,\,v_{p-1},
\]
with $m\ge 0$. Because $G$ is a near-triangulation and $C$ has no
chord, every $u_i$ lies strictly inside $C$, so by hypothesis~(iii)
$|L(u_i)|\ge 5$ for $1\le i\le m$. The graph $G':=G-v_p$ is a
near-triangulation with outer cycle
\[
   C'\colon v_1\,v_2\,\cdots\,v_{p-1}\,u_m\,u_{m-1}\,\cdots\,u_1\,v_1,
\]
on which $v_1,v_2$ remain adjacent and retain their precoloring.

The remainder of the argument is the heart of the proof: we choose
two ``reserve'' colors at $v_p$, modify the lists at the interior
neighbors $u_i$ accordingly, apply the induction to $G'$, and then
return to color $v_p$ from the reserves. To this end, define
\[
   \Omega\;:=\;L(v_p)\setminus\{\sigma(v_1v_p)\,c_1\}.
\]
Since $|L(v_p)|\ge 3$ by hypothesis~(ii), we have $|\Omega|\ge 2$,
and we may choose two distinct values $x,y\in\Omega$. We then
modify the list assignment on $G'$ by setting
\[
   L'(u_i) \;:=\; L(u_i)\setminus\{\sigma(v_pu_i)\,x,\,\sigma(v_pu_i)\,y\}
   \qquad (1\le i\le m),
\]
and leaving $L'(v):=L(v)$ for every other $v\in\V(G')$. Since
$x\ne y$ and $\sigma(v_pu_i)\in\{\pm1\}$, the two excluded values
$\sigma(v_pu_i)\,x$ and $\sigma(v_pu_i)\,y$ are distinct, so at most
two elements of $L(u_i)$ are removed and
\[
   |L'(u_i)|\;\ge\;|L(u_i)|-2\;\ge\;5-2\;=\;3.
\]

We must now verify that the inductive hypotheses apply to
$\Gs|_{G'}$ with the precoloring at $v_1,v_2$ and the list
assignment $L'$. The precoloring of $\{v_1,v_2\}$ is unchanged, so
hypothesis~(i) is inherited directly. For the boundary vertices of
$C'$ other than $v_1$ and $v_2$, two sub-cases arise: vertices in
$\{v_3,\dots,v_{p-1}\}$ retain their original lists, of size at
least $3$ by hypothesis~(ii) of the original problem; while
interior $u_i$, now on $C'$, have $|L'(u_i)|\ge 3$ by the estimate
just established. Either way, hypothesis~(ii) holds. Finally, every
vertex of $G'$ that is interior to $C'$ was already interior to $C$
(no new interior vertices were created), so it retains a list of
size at least $5$, and hypothesis~(iii) holds. By induction, then,
there is a proper $L'$-coloring $c'$ of $\Gs|_{G'}$ extending the
precoloring at $v_1,v_2$.

It remains to extend $c'$ to a coloring $c$ of $G$ by choosing
$c(v_p)\in\{x,y\}$. Such a choice is admissible iff
\[
c(v_p)\;\ne\;\sigma(v_1v_p)\,c_1,\quad
c(v_p)\;\ne\;\sigma(v_{p-1}v_p)\,c'(v_{p-1}),\quad
c(v_p)\;\ne\;\sigma(v_pu_i)\,c'(u_i)\;\;(1\le i\le m).
\]
We check each constraint in turn. First, by the very construction
of $\Omega=L(v_p)\setminus\{\sigma(v_1v_p)\,c_1\}$, both $x$ and
$y$ are different from $\sigma(v_1v_p)\,c_1$, so the constraint
through $v_1$ is automatic. Next, the constraints from the interior
neighbors $u_i$ are eliminated by the modified list assignment: by
construction $c'(u_i)\in L'(u_i)$, so
$c'(u_i)\ne\sigma(v_pu_i)\,x$ and $c'(u_i)\ne\sigma(v_pu_i)\,y$,
and multiplying through by $\sigma(v_pu_i)\in\{\pm1\}$ (using
$\sigma(v_pu_i)^2=1$) gives $\sigma(v_pu_i)\,c'(u_i)\ne x$ and
$\sigma(v_pu_i)\,c'(u_i)\ne y$, so neither $x$ nor $y$ is forbidden
by $u_i$. The only constraint on $c(v_p)$ that may still rule out
an element of $\{x,y\}$ is therefore the one through $v_{p-1}$,
namely $c(v_p)\ne\sigma(v_{p-1}v_p)\,c'(v_{p-1})$, which forbids a
single value. Since $x\ne y$, at least one of $\{x,y\}$ avoids it,
and setting $c(v_p)$ to be such a value yields a proper $L$-coloring
of $\Gs$. This completes the induction.
\end{proof}

\end{document}
